% =============================================================================
% TP2 -- Fundamentos de Inteligência Artificial (UFMG)
% Busca Local e Algoritmos Genéticos no Sequenciamento de Manutenção de Ativos
%
% Compilar (dentro de relatorio/):  pdflatex relatorio && pdflatex relatorio
% Figuras e tabelas em relatorio/figuras/ são geradas, a partir da raiz, por:
%     python scripts/experimentos_t3.py && python scripts/plot_graficos.py
% Figuras ainda inexistentes aparecem como caixas de aviso (o .tex compila mesmo assim).
% =============================================================================
\documentclass[10pt,a4paper,twocolumn]{article}

\usepackage[T1]{fontenc}
\usepackage[utf8]{inputenc}
\usepackage[brazilian]{babel}
\usepackage{newtxtext,newtxmath}            % Times, como no modelo SBC
\usepackage[a4paper,top=2.0cm,bottom=2.0cm,left=1.9cm,right=1.9cm,columnsep=0.55cm]{geometry}
\usepackage{amsmath}
\usepackage{graphicx}
\usepackage{booktabs}
\usepackage{xcolor}
\usepackage{caption}
\usepackage{subcaption}
\usepackage{microtype}
\usepackage[hidelinks]{hyperref}

\captionsetup{font=small,labelfont=bf,labelsep=period}
\setlength{\abovecaptionskip}{3pt}
\setlength{\textfloatsep}{6pt plus 2pt minus 2pt}
\setlength{\dbltextfloatsep}{6pt plus 2pt minus 2pt}
\setlength{\floatsep}{5pt plus 2pt minus 2pt}
\setlength{\dblfloatsep}{5pt plus 2pt minus 2pt}
\makeatletter
\renewcommand\section{\@startsection {section}{1}{\z@}%
  {-1.8ex \@plus -0.5ex \@minus -.2ex}{0.9ex \@plus.2ex}{\normalfont\large\bfseries}}
\renewcommand\subsection{\@startsection{subsection}{2}{\z@}%
  {-1.4ex\@plus -0.4ex \@minus -.2ex}{0.6ex \@plus .2ex}{\normalfont\normalsize\bfseries}}
\renewcommand\paragraph{\@startsection{paragraph}{4}{\z@}%
  {0.9ex \@plus 0.3ex \@minus .2ex}{-0.8em}{\normalfont\normalsize\bfseries}}
\makeatother
\raggedbottom

% Marcador visível para trechos que ainda dependem de outro integrante.
\newcommand{\pendente}[1]{\par\noindent\fcolorbox{orange}{orange!8}{%
  \parbox{\dimexpr\linewidth-2\fboxsep-2\fboxrule}{\small\textbf{[Pendente]} #1}}\par}

% Inclui a figura se o arquivo existir; caso contrário, desenha um aviso.
\newcommand{\figura}[2]{\IfFileExists{#1}{\includegraphics[width=#2]{#1}}{%
  \fbox{\parbox[c][2.5cm][c]{\dimexpr#2-2\fboxsep-2\fboxrule}{\centering\small
  \texttt{\detokenize{#1}}\\ ainda não gerada}}}}
\newcommand{\tabela}[1]{\IfFileExists{#1}{\input{#1}}{\pendente{\texttt{\detokenize{#1}} ainda não gerada.}}}

\title{\vspace{-1.2cm}\textbf{Trabalho Prático 2: Busca Local e Algoritmo Genético}}
\author{%
  \begin{tabular}{@{}c@{\hspace{1.4cm}}c@{\hspace{1.4cm}}c@{}}
    Artur Filardi Victoriano & Bernardo Soares Diniz Lara & Germano Marques Cipriano Fagundes \\
    {\small 2023060707} & {\small 2023060898} & {\small 2023060839}
  \end{tabular}%
}
\date{}

\begin{document}
\maketitle

\begin{abstract}
\noindent
Comparamos três metaheurísticas -- \textit{Greedy Descent} com reinício aleatório (GD),
\textit{Simulated Annealing} (SA) e Algoritmo Genético (GA) -- no problema de sequenciar
$n$ manutenções para minimizar o atraso ponderado total, em instâncias com
$n \in \{12, 50, 150\}$. Todos os métodos recebem o mesmo orçamento de avaliações da
função objetivo e as mesmas sementes aleatórias. Em $n=12$, GD e SA atingem o ótimo
global ($1\,875$) em todas as execuções; com o aumento da escala ($n=50$ e $150$), o SA
passa a superar o GD ($p=5{,}8\times 10^{-4}$ em $n=150$, chegando a superar o valor de
referência do enunciado em $n=50$), enquanto ambos superam significativamente o GA, que
sofre perda rápida de diversidade e baixa eficácia do cruzamento OX diante do
acoplamento temporal dos prazos.
\end{abstract}

% =============================================================================
\section{Formulação do Problema e Validação (T1)}
\label{sec:formulacao}
% Responsável: Integrante 1
Uma equipe parte da base (nó $0$) e deve executar a manutenção de $n$ ativos, um
de cada vez, sem retornar à base ao final. O ativo $j$ fica em $(x_j,y_j)$, exige
$p_j$ unidades de tempo de serviço, tem prazo $d_j$ e peso de criticidade $w_j$. O
deslocamento entre $i$ e $j$ leva
$t_{ij}=\operatorname{round}\!\big(\sqrt{(x_i-x_j)^2+(y_i-y_j)^2}\big)$, com a
matriz pré-calculada. Uma solução é uma permutação $\pi=(\pi_1,\dots,\pi_n)$
dos ativos, com $\pi_0=0$, $C_0=0$ e
\begin{gather}
  C_k = C_{k-1} + t_{\pi_{k-1}\pi_k} + p_{\pi_k}, \notag\\
  h(\pi) = \sum_{k=1}^{n} w_{\pi_k}\max\{0,\; C_k - d_{\pi_k}\},
  \label{eq:h}
\end{gather}
o atraso ponderado total, a ser minimizado. O problema opõe dois critérios.
Visitar primeiro os ativos próximos reduz o tempo gasto em deslocamento, mas pode
deixar para depois um ativo distante, urgente (prazo curto) e crítico (peso alto).
Atender os urgentes primeiro, por sua vez, pode gerar rotas em zigue-zague, que
atrasam todos os ativos seguintes. Como $C_k$ se acumula, uma decisão no início da
rota afeta o atraso de todo o sufixo.

\paragraph{Validação (T1).}
O script \texttt{scripts/tarefa\_t1\_validacao\_h.py} avalia $h$ pela interface usada
pelos algoritmos e por uma implementação direta independente, e reproduz os três
valores de referência do enunciado: $h(1,4,2,3)=39$ e $h(1,4,3,2)=22$ (ótima) no
exemplo da Figura~1, e $h=4\,366$ para a identidade $(1,\dots,12)$ na instância
pequena. O \texttt{round} do Python arredonda empates $x{,}5$ para o par.
Verificamos que nenhuma distância nas três instâncias cai num empate, de modo que
$t_{ij}$ coincide com o arredondamento usual.

% =============================================================================
\section{Algoritmos e Estratégias de Busca}
\label{sec:algoritmos}

\subsection{Vizinhanças e Avaliação Incremental (T2)}
\label{sec:vizinhancas}
% Responsável: Integrante 1
Um movimento é definido por duas posições $i\neq j$ em $\{1,\dots,n\}$. O
\emph{Swap} troca $\pi_i$ e $\pi_j$ e gera $n(n-1)/2$ vizinhos. A \emph{Insertion}
retira $\pi_i$ e o reinsere na posição $j$. Como mover $i$ para $i-1$ equivale a
mover $i-1$ para $i$, ela tem $(n-1)^2$ vizinhos distintos. O \emph{2-opt} inverte o
trecho $\pi_i,\dots,\pi_j$ e gera $n(n-1)/2$ vizinhos. Em $n=150$, são $11\,175$,
$22\,201$ e $11\,175$ vizinhos, respectivamente; um teste automático confere essas
contagens e que os vizinhos são distintos.

\paragraph{Avaliação incremental.}
Um movimento com menor posição alterada $k_{\min}=\min(i,j)$ não altera
$C_1,\dots,C_{k_{\min}-1}$ nem o atraso acumulado
$H_k=\sum_{l\le k} w_{\pi_l}\max\{0,C_l-d_{\pi_l}\}$ desse prefixo. Guardamos
os vetores $C$ e $H$ da solução corrente, e o vizinho é avaliado percorrendo só
o sufixo $k_{\min},\dots,n$, sem alocar vetores. Eles são reconstruídos apenas
quando o movimento é aceito, sem custo de orçamento, pois o $h$ do vizinho já foi
contado. No GD usamos ainda uma \emph{poda}: como $H_k$ é não decrescente, a
varredura para assim que $H_k$ atinge o valor a ser batido. O teste automático
compara a avaliação incremental com a completa em $9\,000$ movimentos aleatórios
(3 vizinhanças $\times$ 3 instâncias) e verifica que a poda nunca descarta uma
melhora.

A Tabela~\ref{tab:t2-speedup} mostra o custo por avaliação, medido sobre toda a
vizinhança de um ótimo local. O ganho é modesto, de $1{,}1$ a $1{,}5\times$. Com
$i,j$ uniformes, $\mathbb{E}[k_{\min}]\approx n/3$, então o sufixo ainda cobre cerca
de $2/3$ da rota, e há um custo fixo por chamada. A poda acrescenta pouco, pelo
próprio formato do problema: no início da rota quase nenhum prazo venceu, e o
atraso se concentra nos últimos ativos. $H_k$ só ultrapassa o limiar perto do
fim, e a varredura raramente é interrompida cedo.

\begin{table}[t]
  \centering\small
  \caption{Custo por avaliação de um vizinho ($\mu$s) num ótimo local: avaliação
  completa, incremental (sufixo) e incremental com poda.}
  \label{tab:t2-speedup}
  \setlength{\tabcolsep}{2.4pt}
  \resizebox{\columnwidth}{!}{\tabela{figuras/tabela_t2_speedup.tex}}
\end{table}

\begin{table}[t]
  \centering\small
  \caption{T2: GD (First Improvement) com reinício por vizinhança, 10 runs com o
  orçamento $B$ do T3. $\tilde h^\star$: mediana do melhor $h$;
  $\bar h_{\text{ót.loc.}}$: $h$ médio dos ótimos locais; ót.\,loc. e passos: por run.}
  \label{tab:t2}
  \setlength{\tabcolsep}{3pt}
  \resizebox{\columnwidth}{!}{\tabela{figuras/tabela_t2_vizinhancas.tex}}
\end{table}

\paragraph{Comparação das vizinhanças (T2).}
A Tabela~\ref{tab:t2} compara as três vizinhanças no GD com reinício. Na
instância pequena, todas atingem o ótimo conhecido $h=1\,875$. Nas maiores, a
Insertion é claramente superior: em $n=150$, a mediana é $80\,011$, contra $90\,544$
(2-opt) e $147\,824$ (Swap), e seus ótimos locais também são os melhores em média. A
explicação está na estrutura do problema. Antecipar um ativo urgente ou adiar um
pouco urgente é exatamente um movimento de inserção, que preserva a ordem
relativa dos demais. O Swap desloca dois ativos ao mesmo tempo, e o 2-opt inverte a
ordem de um bloco inteiro. Essa operação é natural no caixeiro-viajante, mas
aqui reordena prazos que já estavam bem sequenciados. A vizinhança maior da
Insertion faz com que cada descida atinja menos ótimos locais dentro de $B$, mas
melhores. Por isso ela foi adotada como vizinhança padrão do GD e do SA.

\subsection{Greedy Descent com Random Restart}
\label{sec:gd}
% Responsável: Integrante 1
O GD parte de uma permutação aleatória e só aceita movimentos que reduzem $h$.
Implementamos três variantes. O \emph{Best Improvement} avalia toda a vizinhança e
move para o melhor vizinho. O \emph{First Improvement} percorre a vizinhança numa
ordem embaralhada, sorteada a cada reinício e percorrida de forma circular, e
move para o primeiro vizinho que melhora $h$. O \emph{sucessor aleatório} sorteia
vizinhos com reposição e aceita o primeiro que melhora. Nas duas primeiras
variantes, $s$ é um ótimo local quando $|V|$ vizinhos seguidos não melhoram $h$.
Na terceira, o mesmo critério só garante otimalidade local com alta
probabilidade. Ao atingir um ótimo local, o algoritmo o registra e reinicia de uma
nova permutação aleatória, até esgotar $B$. A melhor solução global é mantida à
parte. A solução corrente nunca é reavaliada: seu $h$ é conhecido, e cada
vizinho custa uma avaliação.

\begin{table}[t]
  \centering\small
  \caption{Variantes do GD com Insertion (orçamento e runs do T3). Passos de
  melhoria e ótimos locais por run; aval./passo $=B/$passos;
  $h^\star{=}1875$: runs que atingiram o ótimo conhecido.}
  \label{tab:gd}
  \setlength{\tabcolsep}{3pt}
  \resizebox{\columnwidth}{!}{\tabela{figuras/tabela_gd_variantes.tex}}
\end{table}

\paragraph{Passos dentro do orçamento.}
A Tabela~\ref{tab:gd} mostra o efeito do tamanho da vizinhança. O Best
Improvement paga $|V|$ avaliações por passo: em $n=150$, cada varredura custa
$22\,201$ avaliações, e $B=300\,000$ permite apenas 13 passos de melhoria a
partir de uma solução aleatória. A primeira descida nem termina, e a mediana
final ($822\,160$) é dez vezes pior que a do First Improvement ($80\,011$). O
First e o aleatório dão $\approx 1\,300$ passos com $\approx 200$ avaliações por
passo, porque aceitam a primeira melhora encontrada. O First Improvement foi o
melhor nas instâncias média e grande e é a configuração do GD no T3.

\paragraph{Ótimos locais na instância pequena.}
Em $n=12$, as três variantes atingem o ótimo conhecido $h=1\,875$ em todas as 30
runs. Ao longo dos reinícios, o First Improvement visitou $4\,413$ ótimos locais
(147 por run), mas apenas \textbf{19 permutações distintas}, a mesma contagem do
Best Improvement. A paisagem da Insertion em $n=12$ tem, portanto, poucos ótimos
locais, e $14\%$ das descidas terminam no ótimo global. Com
$\approx 150$ reinícios por run, a chance de nunca cair na bacia do ótimo é
desprezível ($0{,}86^{150}\approx 10^{-10}$). O sucessor aleatório declara 391
``ótimos'' distintos. Pelo critério probabilístico, parte deles não são ótimos
locais de fato, e só $9\%$ das suas descidas terminam em $1\,875$.

% -----------------------------------------------------------------------------
\subsection{Simulated Annealing e Calibração Térmica}
\label{sec:sa}
% Responsável: Integrante 2

\begin{figure*}[t]
  \centering
  \figura{figuras/convergencia.pdf}{\textwidth}
  \caption{Curvas de convergência: mediana do melhor $h$ (escala log) em função das
  avaliações gastas, com faixa interquartil, para as três instâncias. Tracejado:
  mediana do $h$ médio da população do GA.}
  \label{fig:convergencia}
\end{figure*}

O \textit{Simulated Annealing}~\cite{kirkpatrick1983,metropolis1953} é uma busca
local estocástica que escapa de ótimos locais ao aceitar, com probabilidade
controlada, movimentos que pioram a solução. A cada passo sorteia-se uniformemente um par ordenado
(origem, destino) de posições distintas, o que define um vizinho $\pi'$ da solução
corrente $\pi$ na vizinhança Insertion, e calcula-se
$\Delta h = h(\pi') - h(\pi)$. A transição segue o critério de Metropolis, derivado
da distribuição de Gibbs
$P(\pi)\propto e^{-h(\pi)/T}$:
\begin{equation}
  P(\pi \to \pi') =
  \begin{cases}
    1, & \Delta h \le 0,\\[2pt]
    e^{-\Delta h / T}, & \Delta h > 0.
  \end{cases}
  \label{eq:gibbs}
\end{equation}
Como a proposta é simétrica (o movimento inverso tem a mesma probabilidade de
ser sorteado), para $T$ fixo a cadeia de Markov resultante tem a distribuição de
Gibbs como distribuição estacionária. Quando $T\to\infty$, todo vizinho é aceito
(\textit{random walk}). Quando $T\to 0$, a massa de probabilidade se concentra nos
mínimos de $h$, e a regra~\eqref{eq:gibbs} se reduz a uma descida gulosa
estocástica. O resfriamento leva a busca de um regime ao outro. A
melhor solução global $\pi^\star$ é mantida à parte, porque o estado corrente pode
terminar pior do que algum estado já visitado.

\paragraph{Calibração empírica de $T_0$.}
A escala de $\Delta h$ depende fortemente de $n$. Os atrasos se acumulam ao longo
da rota, de modo que reinserir um ativo desloca os instantes $C_k$ de todo o sufixo.
Uma temperatura fixa que funcione em $n=12$ seria fria demais em $n=150$, e uma
adequada a $n=150$ seria um \textit{random walk} em $n=12$. Por isso $T_0$ é
calibrado em cada execução: a partir do estado inicial aleatório, sorteiam-se
200 vizinhos, calcula-se a média $\overline{\Delta h}^{+}$ apenas dos aumentos de
custo ($\Delta h>0$) e impõe-se que, no início, uma piora típica seja aceita com
probabilidade $p_0=0{,}80$:
\begin{equation}
  e^{-\overline{\Delta h}^{+}/T_0} = p_0
  \;\Longrightarrow\;
  T_0 = \frac{-\overline{\Delta h}^{+}}{\ln p_0} \approx 4{,}48\,\overline{\Delta h}^{+}.
  \label{eq:t0}
\end{equation}
A Tabela~\ref{tab:param-sa} mostra que $\overline{\Delta h}^{+}$ cresce cerca de
$15\times$ entre $n=12$ e $n=150$, e $T_0$ acompanha esse crescimento sem nenhum
ajuste manual. As 200 avaliações de calibração são debitadas do orçamento $B$.

\paragraph{Esquema geométrico.}
A temperatura decai a cada avaliação, $T_{k+1}=\alpha\,T_k$, com
\begin{equation}
  T_f = \rho\,T_0,\qquad \alpha = \left(\frac{T_f}{T_0}\right)^{1/L} = \rho^{1/L},
  \label{eq:alpha}
\end{equation}
em que $L = B - 201$ é o número de avaliações que restam após a avaliação do estado
inicial e a calibração. Assim, $T$ atinge $T_f$ exatamente quando o orçamento se esgota.
Adotamos $\rho = 10^{-3}$. Em testes-piloto com 3 a 4 sementes (fora do conjunto
oficial), fixar $T_f$ pela mesma régua de $T_0$ (isto é, $T_f=-\overline{\Delta h}^{+}/\ln p_f$
com $p_f\in[10^{-2},10^{-4}]$) deixava o fim da busca quente demais em $n=150$
(medianas acima de $1{,}5\times10^{5}$). Isso ocorre porque
$\overline{\Delta h}^{+}$ é medido num estado aleatório, cujas pioras são muito maiores
que as observadas perto de bons ótimos locais. Com $\rho=10^{-3}$, $T_f$ fica entre
$2{,}7$ e $39{,}8$ unidades de custo (Tabela~\ref{tab:param-sa}), da ordem das
menores pioras possíveis. Valores menores ($\rho=10^{-4}$) não
trouxeram ganho. A vizinhança é a Insertion, escolhida na T2
(Seção~\ref{sec:vizinhancas}).

\begin{table}[t]
  \centering\small
  \caption{Parâmetros do SA calibrados por instância (mediana entre as runs) e
  qualidade final. $\overline{\Delta h}^{+}$: piora média nos 200 vizinhos iniciais;
  $h^\star_{\min}$ e $\tilde h^\star$: mínimo e mediana do melhor $h$ final.}
  \label{tab:param-sa}
  \setlength{\tabcolsep}{2.6pt}
  \resizebox{\columnwidth}{!}{\tabela{figuras/tabela_parametros_sa.tex}}
\end{table}

\begin{figure}[t]
  \centering
  \figura{figuras/aceitacao_sa.pdf}{\columnwidth}
  \caption{Fração das propostas de piora aceitas pelo SA em janelas de 1\,000
  avaliações, em função da fração consumida do orçamento. Linha: mediana entre as
  runs; faixa: intervalo interquartil.}
  \label{fig:aceitacao}
\end{figure}

\paragraph{Diagnóstico térmico.}
Em cada janela de 1\,000 avaliações registramos a razão entre as pioras aceitas e
as pioras propostas (Figura~\ref{fig:aceitacao}). As três curvas partem de
$\approx 0{,}80$, o que confirma a calibração~\eqref{eq:t0}, e caem de forma suave e
praticamente monotônica: a taxa mediana passa de $0{,}5$ em $\approx 18\%$ do orçamento, de $0{,}2$ em
$\approx 35\%$ e de $0{,}05$ entre $50\%$ e $57\%$, e fica abaixo de $1\%$ a partir de
$\approx 66$--$75\%$. O primeiro terço da busca é, portanto, exploratório:
a solução corrente perambula por regiões de custo alto, o que aparece como a descida
lenta no início das curvas de convergência das instâncias média e grande
(Figura~\ref{fig:convergencia}). O terço final
é, na prática, uma descida gulosa estocástica que refina a bacia escolhida. O
fato mais relevante é que as três curvas quase se sobrepõem quando o eixo
horizontal é normalizado por $B$. A calibração por $\overline{\Delta h}^{+}$
remove a dependência de escala, e a forma do resfriamento passa a depender
apenas da fração do orçamento consumida, não de $n$. Não se observa nenhum dos dois
modos de falha clássicos: nem um platô longo perto de $1$ (temperatura alta demais,
que desperdiça orçamento em passeio aleatório) nem uma queda abrupta para $0$ no
início (têmpera rápida demais, que congela a busca num ótimo local ruim).

\subsection{Algoritmo Genético e Operadores de Permutação}
\label{sec:ga}
% Responsável: Integrante 3
No Algoritmo Genético (\texttt{src/genetic\_algorithm.py}), cada cromossomo é uma
permutação completa $\pi=(\pi_1,\dots,\pi_n)$ precedida pela base $\pi_0=0$. A
população inicial de tamanho $N=100$ é gerada por amostragem uniforme. Como o
problema é de minimização de $h(\pi)$, convertemos o custo na aptidão estritamente
decrescente $f(\pi) = 1/(1+h(\pi))$. A seleção de pais usa \textbf{torneio de
tamanho $k=3$} (garantindo dois pais de índices distintos): sorteiam-se 3 indivíduos
e escolhe-se o de maior $f(\pi)$, o que equivale diretamente a escolher o de menor
custo $h(\pi)$ (o módulo também suporta roleta proporcional a $f(\pi)$). A cada
geração aplicamos \textbf{elitismo} copiando os $E=2$ melhores indivíduos sem
reavaliar $h$, e geramos os $N-E=98$ filhos restantes.

\paragraph{Cruzamento OX, mutação e diversidade.}
Com probabilidade $P_c=0{,}85$, dois pais geram um filho pelo \emph{Order Crossover}
(OX), que preserva a validade da permutação sem duplicar ativos; caso contrário, o
filho copia o primeiro pai. Por exemplo, para $n=6$, $\text{Pai}_1=(1,4,3,5,2,6)$,
$\text{Pai}_2=(6,3,1,2,4,5)$ e corte nas posições $2$ a $4$: (i)~copia-se o bloco
$[4,3,5]$ de $\text{Pai}_1$, formando $(\_,4,3,5,\_,\_)$; (ii)~percorre-se
$\text{Pai}_2$ circularmente a partir da posição $5$, filtrando os ativos ausentes na
ordem em que ocorrem: $[4,5,6,3,1,2]\setminus\{4,3,5\}=[6,1,2]$; (iii)~preenchem-se
as posições $5,6,1$, resultando no filho válido $(2,4,3,5,6,1)$. Em seguida, com
probabilidade $P_m=0{,}20$, aplica-se uma mutação por \emph{Insertion} (vizinhança
escolhida na T2). O orçamento $B$ corresponde a $509$, $1\,529$ e $3\,060$ gerações
completas ($n=12,50,150$). Monitorando a fração de permutações distintas a cada
geração, diagnostica-se convergência prematura: sob torneio $k=3$, a diversidade cai
de $100\%$ para $\approx 20\%$ já em $10\%$ do orçamento ($\approx 50$ gerações) em
$n=12$ e em $50\%$--$60\%$ do orçamento em $n=50$ e $150$, estabilizando em
$18\%$--$24\%$ (exatamente a fração de filhos mutados por geração, $P_m=0{,}20$).

% =============================================================================
\section{Metodologia Experimental e Protocolo T3}
\label{sec:metodologia}
% Responsável: Integrante 2

\paragraph{Orçamento por avaliações, não por tempo.}
O critério de parada de todos os algoritmos é um orçamento $B$ de \emph{avaliações
da função objetivo}. Cada chamada a $h(\cdot)$ consome exatamente uma unidade,
seja ela uma avaliação completa da Eq.~\eqref{eq:h} ou incremental. São contadas a
avaliação da solução inicial, as 200 avaliações de calibração do SA, cada vizinho
examinado pelo GD (inclusive os rejeitados) e cada indivíduo avaliado pelo GA. A
contagem é centralizada em um único objeto \texttt{Avaliador}, encapsulado pela
interface comum \texttt{ProblemaManutencao} (\texttt{src/problema.py}), pela qual
passam todas as chamadas a $h$ e que levanta
uma exceção ao se tentar ultrapassar $B$. Assim, nenhum algoritmo consegue gastar
mais do que o orçamento. O orçamento por avaliações é a métrica de esforço adotada na
comparação porque independe da máquina, da carga do sistema e da qualidade da
implementação. Dois métodos com o mesmo $B$ examinaram o mesmo número de soluções
candidatas. O tempo de parede (\texttt{time.perf\_counter}) é registrado à
parte, como medida secundária. Ele capta diferenças que o orçamento esconde: nem toda
avaliação custa o mesmo (a incremental percorre só as posições $i,\dots,n$), e as
operações genéticas do GA têm custo próprio, que não passa por $h$. No lote T3, o
tempo de parede por avaliação foi de $2{,}6$, $5{,}7$ e $13{,}1\,\mu$s no GD, de
$3{,}0$, $7{,}6$ e $18{,}4\,\mu$s no SA e de $8{,}0$, $14{,}3$ e $31{,}4\,\mu$s no
GA ($n=12$, $50$, $150$). O SA é mais caro por avaliação que o GD porque calcula
$e^{-\Delta h/T}$, sorteia a aceitação e reconstrói $C$ e $H$ a cada movimento aceito,
enquanto o GD só reconstrói ao melhorar e usa poda. Já o GA é o mais caro por
avaliação ($2{,}4\times$ a $3{,}1\times$ o GD) porque todo filho requer avaliação
completa desde a posição $1$, somada ao custo de seleção por torneio, cruzamento OX e
ordenação da elite. Esses tempos são indicativos:
num notebook sob carga contínua, o \emph{throttling} térmico altera o tempo de
parede em até $\approx 30\%$ (uma run isolada do SA em $n=50$ leva $0{,}87$\,s,
contra $1{,}14$\,s de média no lote). Por isso o orçamento em avaliações é a medida
de esforço da comparação.

\paragraph{Configurações.}
As nove combinações (3 algoritmos $\times$ 3 instâncias) usam
$B=50\,000$, $150\,000$ e $300\,000$ avaliações e 30, 30 e 10 runs
independentes nas instâncias pequena ($n=12$), intermediária ($n=50$) e grande
($n=150$), respectivamente. Cada execução parte de uma permutação aleatória e grava, a cada
500 avaliações, o melhor $h$ encontrado até ali, o que dá $B/500$ pontos por curva de
convergência. O SA grava ainda $T$ e a taxa de aceitação de pioras a cada 1\,000
avaliações.

\paragraph{Controle por sementes.}
A run $r$ de qualquer algoritmo usa a semente $s_r = 20\,260\,000 + r$, e toda a
aleatoriedade de uma execução (permutação inicial, sorteio de vizinhos, aceitação,
operadores genéticos) vem de um único gerador \texttt{random.Random}$(s_r)$ local.
Nenhum algoritmo usa o gerador global do módulo \texttt{random}. Com isso,
(i)~qualquer execução é reproduzida bit a bit a partir de sua semente, e
(ii)~os três algoritmos usam as mesmas 30 (ou 10) sementes em cada instância.
Usar a mesma semente, porém, não torna as runs pareadas: cada algoritmo consome
o gerador de forma diferente, de modo que GD, SA e GA na run $r$ não compartilham
solução inicial nem trajetória. As amostras são, portanto, tratadas como
independentes, e as comparações usam o teste de Mann--Whitney
(Seção~\ref{sec:resultados}).

\paragraph{Execução e registro.}
As runs são executadas sequencialmente pelo orquestrador
\texttt{scripts/experimentos\_t3.py}, em um único processo, para que o tempo de
parede não sofra contenção de CPU. Ele grava cada
run em \texttt{resultados\_t3.csv} imediatamente após seu término: algoritmo,
instância, \textit{run}, semente, $h^\star$ final, tempo de parede e histórico de
convergência. Uma interrupção não perde as runs concluídas, e uma nova chamada
retoma das pendentes. Os parâmetros calibrados do SA e o diagnóstico de aceitação vão
para \texttt{parametros\_sa.csv} e \texttt{diagnostico\_sa.csv}. Todas as figuras e
tabelas deste relatório são regeneradas a partir desses arquivos por
\texttt{scripts/plot\_graficos.py}. Os experimentos foram executados em Python~3.11.9,
em um único núcleo de um Intel Core i7-1255U.

% =============================================================================
\section{Resultados Comparativos e Análise Estatística}
\label{sec:resultados}
% Responsável: Integrante 3

\begin{table*}[t]
  \centering\small
  \caption{Resumo do melhor $h$ final por combinação algoritmo--instância, gap
  relativo $(\bar h-1875)/1875$ na instância pequena, orçamento $B$ e tempo médio de
  parede $\bar t$ por run. Valores de referência do enunciado: $1\,875$ ($n=12$),
  $16\,332$ ($n=50$) e $60\,417$ ($n=150$).}
  \label{tab:resultados}
  \tabela{figuras/tabela_resultados_t3.tex}
\end{table*}

\begin{figure}[t]
  \centering
  \figura{figuras/boxplots.pdf}{\columnwidth}
  \caption{Distribuição do melhor $h$ final por algoritmo e instância.}
  \label{fig:boxplots}
\end{figure}

\begin{table}[t]
  \centering\small
  \caption{Experimento adicional do GA (10 runs): cruzamento OX ligado ($P_c=0{,}85$)
  $\times$ desligado ($P_c=0$), $p$-valor de Mann--Whitney, fração de cruzamentos em
  que o filho supera ambos os pais ou ao menos um pai, e diversidade final.}
  \label{tab:ga-exp}
  \setlength{\tabcolsep}{2.5pt}
  \resizebox{\columnwidth}{!}{\tabela{figuras/tabela_ga_experimento.tex}}
\end{table}

A Tabela~\ref{tab:resultados} e os boxplots da Figura~\ref{fig:boxplots} sintetizam as
9 combinações do protocolo T3. Na instância pequena ($n=12$), tanto o GD quanto o SA
encontram o ótimo exato $h^\star=1\,875$ em $30/30$ runs ($\text{gap}=0\%$, $p=1{,}00$
no teste bicaudal de Mann--Whitney). O GA atinge $1\,875$ em apenas $9/30$ runs
($30\%$) e fica retido nos demais ótimos locais da Insertion ($h=1\,937$ em $13/30$
runs e $h=2\,230$ em $7/30$ runs, estagnando em mediana após $8\%$ de $B$), fechando
com mediana $1\,937$, média $1\,999$ ($\text{gap}=6{,}63\%$) e diferença
estatisticamente significativa frente a GD e SA ($p=3{,}7\times 10^{-8}$). Na
instância intermediária ($n=50$), o SA obtém mediana $16\,660$ contra $16\,990$ do GD
($p=0{,}07$, empate estatístico ao nível de $5\%$) e \textbf{alcança mínimo
$h^\star_{\min}=16\,179$, superando o melhor valor de referência do enunciado
($16\,332$)}; ambos superam amplamente o GA (mediana $20\,254$, $p=2{,}4\times
10^{-7}$ para SA$\times$GA e $p=8{,}6\times 10^{-6}$ para GD$\times$GA). Na instância
grande ($n=150$), o SA assume liderança isolada e estatisticamente significativa sobre
o GD (mediana $69\,600$ contra $80\,011$, $p=5{,}8\times 10^{-4}$, com mínimo $66\,238$
próximo à referência $60\,417$), enquanto o GA fica em terceiro lugar ($100\,182$,
$p=1{,}8\times 10^{-4}$ contra SA e $p=4{,}4\times 10^{-4}$ contra GD).

% =============================================================================
\section{Discussão: Escala, Esforço Computacional e Conclusões}
\label{sec:discussao}

\paragraph{Efeito da escala no SA.}
Na instância pequena, o SA atinge o ótimo conhecido $h=1\,875$ em todas as 30 runs, em
mediana após $\approx 21$ mil avaliações ($43\%$ de $B$). O restante do orçamento é
gasto sem melhoria. Nas instâncias média e grande, ao contrário, a curva mediana só
entra a $5\%$ do valor final após $74\%$ e $88\%$ do orçamento, respectivamente, e
a dispersão entre runs cresce: o desvio-padrão relativo vai de $0$ para $4\%$
($n=50$) e $4{,}5\%$ ($n=150$). Isso é coerente com o crescimento do espaço de busca
($n!$) e da vizinhança Insertion ($(n-1)^2 = 22\,201$ vizinhos em $n=150$). Com
$B=300\,000$, o SA examina o equivalente a apenas $\approx 13$ varreduras completas
da vizinhança, e a curva mediana ainda desce quando o orçamento acaba. Isso sugere que, em
$n=150$, o fator limitante é o orçamento, e não o esquema de temperatura.

\paragraph{GD $\times$ SA $\times$ GA: ótimos locais, escala e exploração/intensificação.}
O ranking muda com $n$ (Tabela~\ref{tab:resultados} e
Figura~\ref{fig:convergencia}). Em $n=12$, a paisagem da Insertion possui apenas 19
ótimos locais distintos (Seção~\ref{sec:gd}): o GD (por executar $\approx 147$
reinícios independentes por run) e o SA (por escapar de vales rasos via temperatura)
alcançam o ótimo global $h^\star=1\,875$ em $30/30$ execuções ($6\%$ e $43\%$ de $B$,
respectivamente), enquanto o GA só o atinge em $9/30$ execuções porque sua população
colapsa precocemente sobre um único ótimo local ($1\,937$ ou $2\,230$) sem mecanismo de
reinício. Em $n=50$, o SA termina ligeiramente à frente do GD (mediana
$16\,660$ contra $16\,990$). Em $n=150$, a vantagem do SA é clara: $69\,600$ contra
$80\,011$, uma diferença de $13\%$, com metade da dispersão. As curvas mostram o
mecanismo. No GD, a intensificação é a descida gulosa e a exploração vem só dos
reinícios: em $n=150$, cada descida consome $\approx 230$ avaliações por passo e cabe
pouco mais de um reinício em $B$, extinguindo a exploração por reinício. No
SA, a exploração é a aceitação de pioras controlada por $T$, que escapa de ótimos
locais medianos sem descartar o progresso acumulado; em $n=12$, ao contrário, reiniciar
é a forma mais barata de explorar os 19 ótimos locais.

\paragraph{Quando o GA compensa e experimento sobre recombinação.}
No GA, a exploração cabe à diversidade populacional e à mutação ($P_m=0{,}20$), e a
intensificação cabe à seleção por torneio e ao cruzamento OX. Nas três instâncias, a
curva tracejada da média da população (Figura~\ref{fig:convergencia}) aproxima-se
rapidamente da curva do melhor indivíduo, refletindo a queda da fração de indivíduos
únicos para $\approx 20\%$ (Tabela~\ref{tab:ga-exp}). Nossa hipótese para o desempenho
inferior do GA é que o operador OX, embora preserve a ordem relativa, \textbf{quebra os
instantes acumulados $C_k$}: como o atraso $\max\{0,C_k-d_{\pi_k}\}$ depende da soma de
todos os tempos anteriores, um bloco de ativos bem sequenciado em $\text{Pai}_1$, ao
ser deslocado pelo prefixo herdado de $\text{Pai}_2$, passa a violar prazos. Para
testar essa hipótese diretamente (\texttt{scripts/tarefa\_ga\_experimento.py}),
comparamos o GA padrão ($P_c=0{,}85$) com o mesmo GA sem cruzamento ($P_c=0$, apenas
seleção e mutação). A Tabela~\ref{tab:ga-exp} confirma a hipótese: desligar o OX não
altera significativamente o resultado em nenhuma instância ($p\ge 0{,}52$), e menos de
$1\%$ dos filhos gerados por OX superam ambos os pais. Conclui-se que, sob orçamento
fixo e forte acoplamento temporal no prefixo, buscas locais baseadas na vizinhança
Insertion com avaliação incremental (\textbf{GD} em pequena escala e \textbf{SA} em
média e grande escala) são mais eficazes e até $2{,}8\times$ mais rápidas em tempo de
parede do que o GA puramente geracional.

% =============================================================================
\begin{thebibliography}{9}\small
\bibitem{kirkpatrick1983}
S.~Kirkpatrick, C.~D. Gelatt e M.~P. Vecchi.
\newblock Optimization by simulated annealing.
\newblock \emph{Science}, 220(4598):671--680, 1983.

\bibitem{metropolis1953}
N.~Metropolis, A.~W. Rosenbluth, M.~N. Rosenbluth, A.~H. Teller e E.~Teller.
\newblock Equation of state calculations by fast computing machines.
\newblock \emph{The Journal of Chemical Physics}, 21(6):1087--1092, 1953.

\bibitem{russell2021}
S.~Russell e P.~Norvig.
\newblock \emph{Artificial Intelligence: A Modern Approach}.
\newblock 4.\,ed., Pearson, 2021.
\end{thebibliography}

\end{document}
